% Options for packages loaded elsewhere
\PassOptionsToPackage{unicode}{hyperref}
\PassOptionsToPackage{hyphens}{url}
%
\documentclass[
  10pt,
]{article}
\usepackage{lmodern}
\usepackage{amssymb,amsmath}
\usepackage{ifxetex,ifluatex}
\ifnum 0\ifxetex 1\fi\ifluatex 1\fi=0 % if pdftex
  \usepackage[T1]{fontenc}
  \usepackage[utf8]{inputenc}
  \usepackage{textcomp} % provide euro and other symbols
\else % if luatex or xetex
  \usepackage{unicode-math}
  \defaultfontfeatures{Scale=MatchLowercase}
  \defaultfontfeatures[\rmfamily]{Ligatures=TeX,Scale=1}
\fi
% Use upquote if available, for straight quotes in verbatim environments
\IfFileExists{upquote.sty}{\usepackage{upquote}}{}
\IfFileExists{microtype.sty}{% use microtype if available
  \usepackage[]{microtype}
  \UseMicrotypeSet[protrusion]{basicmath} % disable protrusion for tt fonts
}{}
\makeatletter
\@ifundefined{KOMAClassName}{% if non-KOMA class
  \IfFileExists{parskip.sty}{%
    \usepackage{parskip}
  }{% else
    \setlength{\parindent}{0pt}
    \setlength{\parskip}{6pt plus 2pt minus 1pt}}
}{% if KOMA class
  \KOMAoptions{parskip=half}}
\makeatother
\usepackage{xcolor}
\IfFileExists{xurl.sty}{\usepackage{xurl}}{} % add URL line breaks if available
\IfFileExists{bookmark.sty}{\usepackage{bookmark}}{\usepackage{hyperref}}
\hypersetup{
  hidelinks,
  pdfcreator={LaTeX via pandoc}}
\urlstyle{same} % disable monospaced font for URLs
\usepackage[a4paper,margin=18mm]{geometry}
\setlength{\emergencystretch}{3em} % prevent overfull lines
\providecommand{\tightlist}{%
  \setlength{\itemsep}{0pt}\setlength{\parskip}{0pt}}
\setcounter{secnumdepth}{-\maxdimen} % remove section numbering

\usepackage{esint}
\usepackage{xeCJK}
\setmainfont{DejaVu Serif}
\setmonofont{DejaVu Sans Mono}
\setCJKmainfont{Noto Serif CJK SC}
\setCJKmonofont{Noto Sans Mono CJK SC}
\setmathfont{Asana Math}
\setmathfont{Latin Modern Math}[range=\lrcorner]
\setlength{\emergencystretch}{3em}
\allowdisplaybreaks
\hypersetup{hypertexnames=false}

\author{}
\date{}

\begin{document}

\hypertarget{localization-volume-and-the-first-order-cost-of-whole-fillings}{%
\section{Localization volume and the first-order cost of whole
fillings}\label{localization-volume-and-the-first-order-cost-of-whole-fillings}}

28 September 2026. Continuation from \texttt{105df84}, using the
relative identities of \href{20260928_relative_cell_cost.md}{I54}. The
estimates below apply to complete smooth fillings, not selected terminal
bubbles. They settle the first coefficient for a localized class with
arbitrary interior metrics. They do not extract that class from
arbitrary near-minimizers, so O11 remains open.

\hypertarget{an-exact-all-metric-bound-with-a-fixed-exterior}{%
\subsection{1. An exact all-metric bound with a fixed
exterior}\label{an-exact-all-metric-bound-with-a-fixed-exterior}}

Let X be a closed complex-hyperbolic surface, with its complex
orientation and metric satisfying Ric = -g. Write A = c1 squared(X), so
its volume is 2 pi squared A. Remove finitely many disjoint smoothly
bounded regions which are balls, with total background volume D. Replace
them by complete marked smooth cells of the topological types in I54,
Section 1. The metric agrees with the original CH metric on the full
outer collars and on the entire complement. There are no hypotheses on
the interior metric, concentration, diameter, volume, derivatives, or
its distance from a model.

Let k and ell be the total blowup and handle numbers and put

\[
 Q=k+4\ell>0,\qquad a=2\pi^2A-D>0,\qquad
 b=8\pi^2Q-4D,\qquad c=4a+b.
\]

In this section assume b is nonnegative. On the union C of the inserted
cells define

\[
 v=\operatorname{Vol}(C),\quad t=\int_Cs,\quad
 S=\int_Cs^2,\quad W=\int_C|W^+|^2.
\]

The relative identities sum over whole cells; no interior boundary is
introduced. The total Ricci energy R and total scalar integral I are

\[
 R=c+S/3+4W,\qquad I=-4a+t,\qquad V=a+v.                 \tag{1.1}
\]

In particular R is positive. Recall the scale-invariant residual

\[
 \mathcal E=\inf_{\lambda\ne0}
 \frac{\int|\operatorname{Ric}-\lambda g|^2}{\lambda^2 V}
          =4-\frac{I^2}{VR}.
\]

The formula remains valid when I = 0, with infimum four as the absolute
value of lambda tends to infinity. No normalization of the unknown
optimum lambda is imposed.

\textbf{Theorem 1.} Every metric with this complete exterior satisfies

\[
 \boxed{\mathcal E\ge
 \frac{4ab/c+v}{a+v}
 =\frac{4b}{c}+\left(1-\frac{4b}{c}\right)\frac{v}{a+v}
 \ge\min\{4b/c,1\}.}                                   \tag{1.2}
\]

For b at most 4a/3, the last lower bound is explicitly

\[
 \boxed{\mathcal E\ge
 \frac{4(Q-D/(2\pi^2))}{A+Q-D/\pi^2}.}                  \tag{1.3}
\]

\emph{Proof.} Cauchy gives S at least t squared/v. Applying weighted
Cauchy a second time yields

\[
 (-4a+t)^2\le
 \left(\frac{16a^2}{c}+3v\right)
 \left(c+\frac{t^2}{3v}\right)
 \le\left(\frac{16a^2}{c}+3v\right)R.
\]

Subtract the resulting upper bound for I squared/(VR) from four. This
proves (1.2); inserting a and b proves (1.3). These uses of Cauchy are
on actual integrals. No formally optimized moments are claimed to be
realized by a smooth metric. The underlying relative curvature identity
is classical Chern--Weil, as derived in I54, Section 1. QED.

The bound does not require an actual monopole on an individual cell. The
original closed family's actual supply is preserved because the
topological surgery types and their markings have not changed. In
particular this is an application to that family, not a construction of
a new supplier.

\hypertarget{the-equality-remainder-and-the-whole-core-energy-distribution}{%
\subsection{2. The equality remainder and the whole-core energy
distribution}\label{the-equality-remainder-and-the-whole-core-energy-distribution}}

To retain the equality information, set

\[
 \delta=\frac13(S-t^2/v)+4W\ge0,\qquad
 K=16a^2/c+3v,\qquad H=\frac{(3vc+4at)^2}{3vc}.
\]

Since R = c + t squared/(3v) + delta, exact expansion gives

\[
 KR-I^2=K\delta+H.
\]

Thus the proof of Theorem 1 retains the identity

\[
 \boxed{\mathcal E=
 \frac{4ab/c+v}{a+v}
       +\frac{K\delta+H}{(a+v)R}.}                       \tag{2.1}
\]

Consider the original balanced sequence with Q\_j = A\_j h\_j, h\_j
tending to zero, and assume only that D\_j = o(Q\_j). The following
result also applies to other surgery proportions on CH covers; the
balanced proportions are not used in its proof.

\textbf{Theorem 2.} In this class, independently of the interior
geometry,

\[
 \liminf_j\mathcal E(g_j)/h_j\ge4.                       \tag{2.2}
\]

Furthermore the following two statements are equivalent:

\[
 \mathcal E(g_j)/h_j\longrightarrow4;
 \qquad v_j+S_j+W_j=o(Q_j).                              \tag{2.3}
\]

The three summands in the second statement are nonnegative.
Equivalently, each is o(Q\_j). For such a sequence,

\[
 \int_{C_j}|\operatorname{Ric}|^2=8\pi^2Q_j+o(Q_j),
 \qquad \int_{C_j}|\operatorname{Ric}^{\circ}|^2
                                  =8\pi^2Q_j+o(Q_j),    \tag{2.4}
\] \[
 \int_{C_j}|W^-|^2=12\pi^2k_j+o(Q_j).                   \tag{2.5}
\]

These are statements about the entire inserted region, including all
intermediate necks and scales, not a classification of individual
bubbles.

\emph{Proof.} As j tends to infinity, a is asymptotic to 2 pi squared A,
c is asymptotic to 4a, and L := 4b/c = 4h + o(h). Eventually L is less
than one. Equation (1.2) first proves (2.2). If E/h tends to four, it
also gives v/(a+v) = o(h), hence v = o(Q).

One must use the remainder, not infer small scalar energy from this
volume statement alone. Since t squared is at most vS and S is at most
3R, the exact definition of E gives

\[
 \sqrt{(4-\mathcal E)(a+v)}\sqrt R
       =|-4a+t|\le4a+\sqrt{3v}\sqrt R.
\]

Here E tends to zero and v/a tends to zero. Moving the last term to the
left proves R = O(a); also R is at least c, which is comparable to a.
The first term of (2.1) is 4h + o(h), so

\[
 K\delta+H=o(Qa).
\]

Because K is comparable to a, delta = o(Q). Also

\[
 H=\frac{16a^2}{3vc}\left(t+\frac{3vc}{4a}\right)^2,
 \qquad
 \left(t+\frac{3vc}{4a}\right)^2=o(Qv).
\]

As c/a is bounded and v = o(Q), it follows that t squared/v = o(Q). The
definition of delta now proves S + W = o(Q).

Conversely, v + S + W = o(Q) gives t = o(Q) by Cauchy. In (1.1), R = 4a
+ b + o(Q), I = -4a + o(Q), V = a + o(Q). Expanding the exact expression
for E gives E = 4h + o(h). Finally the relative Ricci and signature
identities in I54 give (2.4)--(2.5). QED.

Classical shrinking surgery constructions attain (2.3) by choosing their
scales sufficiently small. In this localized class the lower coefficient
is therefore sharp. The theorem does not assume those constructions
describe all competing interiors.

\hypertarget{a-fixed-improvement-must-modify-a-nonnegligible-amount-of-background}{%
\subsection{3. A fixed improvement must modify a nonnegligible amount of
background}\label{a-fixed-improvement-must-modify-a-nonnegligible-amount-of-background}}

Let theta = D/(2 pi squared Q). When 0 \textless= theta \textless{} 1
and h is small, Theorem 1 gives

\[
 \frac{\mathcal E}{h}\ge
        \frac{4(1-\theta)}{1+h-2\theta h}.                \tag{3.1}
\]

If e := E/h is at most 4-epsilon, with fixed positive epsilon and small
h, rearrangement gives

\[
 \theta\ge\frac{4-e(1+h)}{4-2eh}.
\]

The branch theta \textgreater= 1 already has a stronger volume lower
bound. Consequently any such subfour sequence in the complete-exterior
class must satisfy

\[
 \boxed{\liminf_j\frac{D_j}{Q_j}
                       \ge\frac{\pi^2}{2}\varepsilon.}  \tag{3.2}
\]

It cannot obtain a fixed first-order improvement by arbitrary violent
changes confined to background regions of total volume o(Ah). This is
different from the shrinking C0-neighborhood theorem I53: the interior
here need not be close to any recovery metric, but the entire exterior
is fixed. Neither hypothesis is asserted to follow from the other.

For a fixed cell replicated sparsely, I54's stronger supplier-based
bound J \textgreater= 8 pi squared q - 4d/3 gives the sharper necessary
cost condition d/q \textgreater= (3 pi squared/2) epsilon for a
coefficient at most 4-epsilon. The all-metric bound above does not
presuppose fixed cells, uniform cell volume, or convergence of their
individual energies.

The unresolved alternative is now specific: a lower coefficient must
change a Q-scale amount of background, or must not admit this marked
complete-exterior description. No extraction theorem has been proved.

\hypertarget{individual-cores-and-their-topology-need-not-be-identified}{%
\subsubsection{3.1 Individual cores and their topology need not be
identified}\label{individual-cores-and-their-topology-need-not-be-identified}}

The preceding proof extends further than its cell formulation. Suppose
the closed manifold M is already the prescribed member of the actual
family, hence \(2\chi(M)+3\tau(M)=A-Q\) and \(\tau(M)=\tau(X)-k\).
Assume only that an open subset U of the CH base X is
orientation-preservingly isometric to an open subset U' of M. Put D =
Vol(X without U) and let C be M without U'.

Subtract the two \emph{closed} Chern--Weil identities and then split
their integrals into the common open set and its complement. The same
equations (1.1) and the relative signature identity follow. No Stokes
theorem on the complement and no boundary transgression are needed.
Theorems 1--2 and (3.2) are unchanged. The complements may be clustered,
connected, multiply bounded, or lack a smooth boundary; none has to be
identified as a ball or an individual surgery core. All integrals over
them are ordinary measurable-set integrals of the smooth closed metric.

Thus the statement does not require a bubble-tree decomposition. It only
requires the common metric region and the known topology of the closed
member of the family. The common-region hypothesis itself has not been
derived from small Einstein defect.

In fact no identification with a region in a specified base is needed.
It is enough that M, of the same known closed topology, has a complex-
oriented locally CH open region of volume at least 2 pi squared A -
o(Q), normalized by Ric = -g. Select a measurable subregion with volume
at most 2 pi squared A and discrepancy o(Q). Its scalar, Ricci and Weyl
densities are exactly those used in (1.1). Splitting the closed
identities gives (2.2) directly, without an external marking or any
asserted topology of its complement.

\hypertarget{the-exterior-need-not-be-einstein}{%
\subsection{4. The exterior need not be
Einstein}\label{the-exterior-need-not-be-einstein}}

There is an exact interface for a varying exterior. Start with any
smooth reference metric on the closed base, remove the reference cells
D, and replace them by the same marked surgeries, preserving the full
metric germs. On the unchanged exterior let

\[
 a=\operatorname{Vol}(M\setminus C),\quad
 i=\int_{M\setminus C}s,\quad
 r=\int_{M\setminus C}|\operatorname{Ric}|^2>0,
 \quad e_{\mathrm{out}}=4-i^2/(ar)\ge0.
\]

All reference cap integrals below concern the reference metric, not the
inserted geometry. Relative Chern--Weil gives

\[
 b=8\pi^2Q+R_D-S_D/3-4W_D^+,\quad c=r+b,\quad
 R=c+S/3+4W,\quad I=i+t.                                 \tag{4.1}
\]

Assume b is nonnegative. Set

\[
 L=4-\frac{i^2}{ac}
  =e_{\mathrm{out}}+(4-e_{\mathrm{out}})\frac b{r+b},
 \quad K=i^2/c+3v,\quad H=\frac{(3vc-it)^2}{3vc}.
\]

Exactly as before, with delta = (S-t squared/v)/3 + 4W,

\[
 \boxed{\mathcal E=\frac{aL+v}{a+v}
                 +\frac{K\delta+H}{(a+v)R}.}             \tag{4.2}
\]

This formula is valid even when i vanishes. It explicitly retains the
cost of a non-Einstein exterior; replacing the cap correction by the CH
value -4D without its hypotheses would be incorrect.

\textbf{Theorem 3.} Suppose h = Q/A tends to zero,

\[
 r=8\pi^2A+o(A),\qquad b=8\pi^2Q+o(Q),
                    \qquad\mathcal E=O(h).              \tag{4.3}
\]

Then e\_out = O(h), v/a = O(h), and

\[
 \boxed{\mathcal E\ge e_{\mathrm{out}}+4h+\frac{v}{a+v}+o(h).}
                                                               \tag{4.4}
\]

Moreover E/h tends to four if and only if

\[
 e_{\mathrm{out}}=o(h),\qquad v/a=o(h),\qquad S+W=o(Q).
                                                               \tag{4.5}
\]

The condition E = O(h) in (4.3) is dispensable in the reverse
implication of (4.5), since the other assumptions then give it.

\emph{Proof.} Since E tends to zero and (4.2) implies E \textgreater=
min(L,1), eventually L \textless= E. Also L \textgreater= e\_out, so
e\_out = O(h). The first term of (4.2) is L + (1-L)v/(a+v); hence v/a =
O(h). Assumptions (4.3) then give L = e\_out + 4h + o(h), proving (4.4).

If E/h tends to four, (4.4) forces e\_out = o(h) and v/a = o(h). As in
Section 2,

\[
 \sqrt{(4-\mathcal E)(a+v)}\sqrt R
                 \le |i|+\sqrt{3v}\sqrt R
                 \le2\sqrt{ar}+\sqrt{3v}\sqrt R.
\]

Thus R = O(r). Now K is comparable to a, c is comparable to r, and the
remainder in (4.2) is o(h). It follows that delta = o(Q), H = o(aQ), and

\[
 \left(t-\frac{3vc}{i}\right)^2=o(vQ).
\]

Here i is nonzero for large j since e\_out tends to zero. Also
\(v c^2/i^2=O(vr/a)=o(Q)\). Consequently t squared/v = o(Q), and S+W =
o(Q). Conversely these conditions give R = r + 8 pi squared Q + o(Q).
The scalar cross term obeys

\[
 \frac{|it|}{ar}\le
             2\sqrt{\frac{v}{a}\frac{S}{r}}=o(h).
\]

The same exact formula for E proves the reverse implication. QED.

For example, (4.3)'s cap condition holds if R\_D + S\_D + W\_D\^{}+ =
o(Q); this is a property of actual reference caps, not an assumed bubble
decomposition. The exterior energy condition must likewise be supplied.
Theorem 3 gives additive first-order costs of bulk defect and topology
in this precise interface. It is not a Gamma-convergence theorem and
gives neither compactness nor an automatic choice of reference metric
for every near-minimizer.

Here is a natural way to supply the reference energy in (4.3). Suppose
the reference metrics on the CH bases themselves obey
\(\mathcal E(g_j^{\mathrm{ref}})=O(h_j)\), and the removed reference
caps have \(R_D+W_D^+=o(Q_j)\). The classical Ricci-energy lower bound
on a complex-hyperbolic surface is \(R_{\mathrm{ref}}\ge8\pi^2A\); see
\href{https://arxiv.org/html/math/0003068}{LeBrun, Theorem 4.3}. The
opposite inequality needed here follows directly from the signature:

\[
 W_{\mathrm{ref}}^+\ge12\pi^2\tau(X)=4\pi^2A/3,\qquad
 R_{\mathrm{ref}}=4P_{\mathrm{ref}}-12W_{\mathrm{ref}}^+
                                                    +24\pi^2A,
\] \[
 4P_{\mathrm{ref}}\le
       \mathcal E(g^{\mathrm{ref}})R_{\mathrm{ref}},
 \qquad
 R_{\mathrm{ref}}\le
       \frac{8\pi^2A}{1-\mathcal E(g^{\mathrm{ref}})}.    \tag{4.6}
\]

For the last estimate take j large enough that the denominator is
positive. Since S\_D \textless= 4R\_D, these hypotheses give exactly b =
8 pi squared Q + o(Q) and r = 8 pi squared A + O(Q). Thus Theorem 3
applies to arbitrary replacement interiors on these non-Einstein
reference backgrounds. No smallness of their C0 distance from the CH
metric is required. The cap smallness remains an actual assumption, not
a consequence of a scalar or Ricci bound alone.

\hypertarget{an-intrinsic-bulk-criterion-without-a-reference-metric}{%
\subsubsection{4.1 An intrinsic bulk criterion without a reference
metric}\label{an-intrinsic-bulk-criterion-without-a-reference-metric}}

The complete region-matching assumption can be replaced by precisely the
integral information used in the proof. On any closed member M with
\(T_+=A-Q\), choose any measurable subset U of positive volume, with
complement C of positive volume. Define its curvature budget

\[
 T_U=\frac1{8\pi^2}\left\{\frac13\int_U s^2
                   +4\int_U|W^+|^2-\int_U|\operatorname{Ric}|^2\right\}.
                                                               \tag{4.7}
\]

This is a real-valued curvature integral, not an integer or a relative
characteristic number. The closed Chern--Weil identity gives (4.1) with
\[
                  b=8\pi^2(T_U-T_+).
\] Consequently, if actual sets U\_j can be chosen with \[
 T_{U_j}=A_j+o(Q_j),\qquad
 \int_{U_j}|\operatorname{Ric}|^2=8\pi^2A_j+o(A_j),        \tag{4.8}
\] Theorem 3 applies verbatim. Neither a metric model on U\_j, a smooth
interface, a relative gauge marking, nor an identification of the
complement's topology is required.

Its fine additive and equality statements still concern competitive
sequences with E = O(h). Conditions (4.8) alone imply the liminf lower
coefficient four, not the upper bound E = O(h).

This is a sufficient integral selection criterion, not a completed
selection theorem. In particular, the global identity only gives
\(T_M=A-Q\); it does not supply a subset carrying \(A+o(Q)\). Equation
(4.8) expresses exactly the missing first-order separation of the bulk
curvature budget from the negative budget in its complement. If the
actual field equations produce such sets for all competitive sequences,
the all-metric lower coefficient follows. The present proof does not
assume that conclusion.

\hypertarget{a-finite-amplitude-exclusion-for-the-coupled-diagonal-profiles}{%
\subsection{5. A finite-amplitude exclusion for the coupled diagonal
profiles}\label{a-finite-amplitude-exclusion-for-the-coupled-diagonal-profiles}}

For the remaining fixed-scale construction problem, use precisely the
actual profiles F,w,phi and the symbols D\_0,D\_3,D\_H,c\_0,z\_0,U of
I54, Section 6.3. Subscripts on c\_0,z\_0 in this section distinguish
them from the global moment c above. Write

\[
 G_w=e^{2\phi}K_{F,w},\quad G_0=e^{2\phi}K_{F,0},
 \quad V_w=U-2F\phi'w'.
\]

The full fixed collar is CH; the ordinary bolt is at r\_0 = sqrt(mu). In
particular F(r\_0) = 0, F'(r\_0) = 2/r\_0, and w has the smooth
ordinary-bolt extension specified in I54. Direct subtraction of the
complete cost, including every transition, gives

\[
\begin{aligned}
 \frac{J(G_w)-J(G_0)}{2\pi^2}=\int r^3\bigl[
 &-4D_0F(w')^2+4F^2(w')^4+4D_Hc_0+2c_0^2\\
 &+2V_w^2-2D_3z_0+z_0^2\bigr]\,dr.
\end{aligned}\tag{5.1}
\]

Fix r\_c,R,K,M independent of mu, with r\_0 \textless{} r\_c \textless{}
R and R inside the unchanged CH collar. On the fixed interval
{[}r\_c,R{]}, assume

\[
 M^{-1}\le F\le M,\quad |F'|+|\phi'|\le M,\quad |w|\le K,\quad
                   \max(|D_0|,|D_3|,|D_H|)\le\eta.
                                                               \tag{5.2}
\]

\textbf{Proposition 4.} There is a positive eta\_0 depending only on
r\_c,R,K,M such that, if eta \textless= eta\_0 and

\[
 D_0\le0,\qquad D_3\le0,\qquad D_H\ge0
                                     \quad(r_0<r<r_c),  \tag{5.3}
\]

then J(G\_w) \textgreater= J(G\_0) \textgreater{} 8 pi squared. The
inner anisotropy has no amplitude or derivative bound, apart from
smoothness.

\emph{Proof.} On the outer interval, the definition of V\_w becomes

\[
 w''+\left(F'/F+3/r+2\phi'\right)w'+b_0(r,w)w=-V_w/F.
\]

Here b\_0 extends smoothly across w = 0 and is uniformly bounded by
(5.2); the quotient sinh(aw)/w is bounded by a cosh(aK). The terminal
data w = w' = 0 are supplied by the full collar. The elementary
terminal-value estimate for this bounded first-order system gives

\[
 \int_{r_c}^{R}r^3(w^2+(w')^2)\,dr
                      \le C\int_{r_c}^{R}r^3V_w^2\,dr.  \tag{5.4}
\]

One can apply Gronwall to the coefficients evaluated at the actual
solution w; only their bounds are used. Since c\_0+z\_0 \textless= Cw
squared on this fixed interval, the outer part of (5.1) is at least 2
integral V\_w squared minus C eta integral(w squared+w' squared). Choose
eta\_0 so that (5.4) absorbs the negative term. All terms of (5.1) on
the inner interval are nonnegative by (5.3). Finally G\_0 is conformal
to the genuine Kahler ordinary blowup K\_\{F,0\}. I54's whole-cell
conformal bound, with exceptional area pi mu, gives

\[
 J(G_w)\ge J(G_0)\ge
       \frac{2\pi^2d(\mu+2)^2}{d-\pi^2\mu^2/2}
       +\frac14\int_C(s_{G_0}+4)^2
                    >8\pi^2.                           \tag{5.5}
\]

The denominator is the positive volume of the canonical Kahler cell. In
particular the bound is at least 8 pi squared + 8 pi squared mu + 2 pi
squared mu squared. QED.

These hypotheses hold on a uniform, explicitly weighted profile
neighborhood of the low-cost recovery. On the inner region use the uncut
cscK profile F\_mu of I54 and \(p_\mu=\mu(\mu+2)/r^4\), and write F =
F\_mu+f.~It suffices that

\[
\begin{aligned}
 |f''|+|f'|/r+|f|/r^2+|e^{2\phi}-1|
 +F|\phi''|+(|F'|+F/r)|\phi'|+F(\phi')^2
                                      \le p_\mu/10.
\end{aligned}\tag{5.6}
\]

The definitions give \textbar D\_0+p\_mu\textbar,
\textbar D\_3+p\_mu\textbar, \textbar D\_H-p\_mu\textbar{} \textless=
p\_mu/2. On the fixed outer interval a sufficiently small ordinary C2
neighborhood of the complete recovery, followed by sufficiently small
mu, gives (5.2). The radii and thresholds do not depend on mu. This is a
weighted second-derivative neighborhood, not a claim that low energy
automatically supplies those bounds.

If only (5.2) is known, keep the potentially negative inner payment

\[
 \mathfrak C_\mu=\int_{r_0}^{r_c}r^3
 \{4(D_0)_+F(w')^2+4(-D_H)_+c_0+2(D_3)_+z_0\}\,dr.
\]

The same proof gives J(G\_w) \textgreater= J(G\_0)-2 pi squared C\_mu.
Thus an actual subthreshold cell in this outer regime necessarily has

\[
 \boxed{\mathfrak C_\mu>
       \frac{d(\mu+2)^2}{d-\pi^2\mu^2/2}-4
                                      \ge4\mu+\mu^2.}   \tag{5.7}
\]

The sign violations must correlate with actual anisotropy strongly
enough to pay this amount. Their mere occurrence is not a construction.
No bound of the reverse kind in (5.7) has been proved for arbitrary
profiles, and (5.2) has not been extracted from all low-energy cells.

\hypertarget{the-actual-correlation-in-an-almost-kahler-filling}{%
\subsection{6. The actual correlation in an almost-Kahler
filling}\label{the-actual-correlation-in-an-almost-kahler-filling}}

For the nonconformal metric K itself, there is a geometric version of
the missing signed payment. Use the ordinary blowup, positive periods
t\_i and relative first-Chern class of I54, Section 6.1. Locally choose
a J-anti-invariant two-form psi of squared norm two and write \[
 \nabla\omega=\zeta\otimes\psi-J\zeta\otimes J\psi,\qquad
 n=\tfrac12|\nabla\omega|^2=2|\zeta|^2.
\] The classical canonical Chern form is \[
 \gamma=\operatorname{Rm}(\omega)-\zeta\wedge J\zeta,\qquad
 [\gamma]=2\pi c_1.
\] These are formulas (21) and (25) in
\href{https://faculty.fiu.edu/~draghici/research/paperfiles/Math-Ann.pdf}{Apostolov--Armstrong--Draghici}.
Put \[
 \xi=(\zeta\wedge J\zeta)^-,\quad |\xi|^2=n^2/8,\quad
 \beta=W^+(\omega)-(s/6+n/2)\omega,
\] \[
 \rho'_0=(\operatorname{Ric}'-sK/4)(J\cdot,\cdot).
\] Primes denote the J-invariant part; double primes denote the
J-anti-invariant part. In particular
\(|\operatorname{Ric}'_0|^2=2|\rho'_0|^2\). The actual form satisfies \[
 (\gamma+\omega)^+=(s+n+4)\omega/4+\beta,\qquad
 (\gamma+\omega)^-=\rho'_0-\xi.
\] On the full CH collar gamma + omega vanishes. After gluing the
unchanged closed CH exterior, its cohomology class has only exceptional
components. Thus, with \(C_E=\sum_i(2\pi+t_i)^2\), \[
                    \int_C(\gamma+\omega)^2=-C_E.
\] Expanding this identity and using
\(J(K)=2\int|\rho'_0|^2+\int|\operatorname{Ric}''|^2 +\tfrac14\int(s+4)^2\)
gives \[
 \boxed{\begin{aligned}
 J(K)={}&2C_E+\tfrac12\int(s+4)^2+\int\mathcal C\\
        &+2\int|\beta|^2+\int|\operatorname{Ric}''|^2,
 \end{aligned}}                                                   \tag{6.1}
\] where the signed smooth density is \[
 \mathcal C=2(\operatorname{Ric}+K)(\zeta^\#,\zeta^\#)
       +2(\operatorname{Ric}+K)((J\zeta)^\#,(J\zeta)^\#).
\] Away from n = 0 it is
\(n\,\operatorname{tr}_{D^\perp}(\operatorname{Ric}+K)\), where
\(D^\perp\) is the two-plane spanned by zeta and J zeta. The displayed
smooth formula does not require extending that plane across the zero
set. Rotating psi rotates this pair, leaving the density unchanged.
Algebraically the term is \(n(s+4)/2+4\langle\rho'_0,\xi\rangle\).

Consequently a genuinely subthreshold AK cell must have a sufficiently
negative torsion-weighted directional Ricci correlation. The integrated
nonintegrability N alone does not decide its sign. Nor does a small
exceptional period force N to vanish: a fixed compactly supported
compatible nonintegrable perturbation away from the shrinking bolt
retains positive N in the limit. Such an example pays the positive cost
of the background perturbation and is not subthreshold. Equation (6.1)
is a consequence of the classical canonical Chern identity, not a new
integrability theorem. Its analogue for arbitrary conformal metrics
retains the conformal terms in I54 rather than silently identifying G
with K.

\hypertarget{full-coupled-hessian-check-at-the-low-cost-recovery}{%
\subsection{7. Full coupled Hessian check at the low-cost
recovery}\label{full-coupled-hessian-check-at-the-low-cost-recovery}}

The nonlinear result in Section 5 does not depend on a formal Taylor
expansion. The following separate calculation records why a quadratic
search near the entire cutoff recovery does not produce a missing
negative mode.

Choose fixed physical radii \(r_c<r_1<R\), cut the cscK profile F\_mu to
CH between r\_c and r\_1, and write \(s_{F_\mu}=-4+\delta_\mu\). Then
delta\_mu is supported in that annulus with C1 norm O(mu). Consider
\(F=F_\mu+\epsilon f,\ \phi=\epsilon p,\ w=\epsilon z\), with the actual
fixed-bolt and collar conditions; in particular f(r\_0) = f'(r\_0) = 0
and all variations vanish near R. Put \[
 \mathcal Sf=-f''-7f'/r-8f/r^2,\quad
 L_Fp=Fp''+(F'+3F/r)p'.
\] For \(j=J/(2\pi^2)\), the even quadratic coefficient is \[
 Q_{\mathrm{ev}}=\mathcal Q_{\mathrm{ev}}
 +\int r^3\{4\delta_\mu'fp'-4\delta_\mu F_\mu(p')^2
                                      +4\delta_\mu p^2\}\,dr,
\] \[
 \mathcal Q_{\mathrm{ev}}=\int r^3
 \left\{\frac{(\mathcal Sf)^2}{6}
 +\frac{(\mathcal Sf-6L_{F_\mu}p+6p)^2}{3}
                            +4p^2+4F_\mu(p')^2\right\}\,dr.
                                                               \tag{7.1}
\] This follows by expanding the exact even-sector expression \[
 j=4-\frac{2d}{\pi^2}+\int r^3
       \{s_F^2/6+X^2/12+(X+4e^{2\phi})^2/4\}\,dr,
 \quad X=s_F-6L_F\phi-6F(\phi')^2.
\] The integration by parts uses \(\int r^3L_fp=[r^3fp']_{r_0}^{R}=0\).
The terminal ODE for S f on the fixed outer annulus controls its H1 norm
by the L2 norm of S f there. Consequently
\(Q_{\mathrm{ev}}\ge(1-C\mu)\mathcal Q_{\mathrm{ev}}\).

For the odd sector define \[
 D_Fz=-Fz''-(F'+3F/r)z'+(16/(r^2F)-8/r^2)z.
\] Its coefficient is \[
 Q_{\mathrm{odd}}=\int r^3
 \{2(D_{F_\mu}z)^2-4E_VF_\mu(z')^2
                     +32E_Hz^2/r^2-64E_Vz^2/(r^2F_\mu)\}\,dr.
\] In the cscK core \(E_V=-p_\mu,E_H=p_\mu\). In the fixed transition
annulus these two defects are O(mu). The terminal ODE estimate, now for
D\_F z, therefore gives \[
\begin{aligned}
 Q_{\mathrm{odd}}\ge{}&(2-C\mu)\int r^3(D_{F_\mu}z)^2\,dr\\
 &+\int_{r_0}^{r_c}r^3p_\mu
       \{4F_\mu(z')^2+32z^2/r^2+64z^2/(r^2F_\mu)\}\,dr.
\end{aligned}                                                     \tag{7.2}
\] Mixed even--odd second derivatives vanish because the full functional
is even in w. Thus the coupled quadratic form is positive for small mu.
The linear derivative is \[
 Dj[f,p,z]=\int r^3\delta_\mu
                           (\mathcal Sf-4L_{F_\mu}p+4p)\,dr.
\] It is bounded by \(C\mu\sqrt{\mathcal Q_{\mathrm{ev}}}\); quadratic
optimization can save only O(mu squared), whereas \[
 j(F_\mu,0,0)=(\mu+2)^2+\tfrac12\int r^3\delta_\mu^2\,dr
                                      =4+4\mu+O(\mu^2).
\] This is a statement about the quadratic model. Finite deformations
are justified only under Section 5's actual nonlinear hypotheses. The
search tool fixes its cutoff in logarithmic coordinates instead; its
moving physical transition does not automatically inherit these
fixed-annulus uniform constants.

\hypertarget{a-global-radial-moment-the-entire-nonconformal-triaxial-class}{%
\subsection{8. A global radial moment: the entire nonconformal triaxial
class}\label{a-global-radial-moment-the-entire-nonconformal-triaxial-class}}

There is a stronger exclusion than the neighborhood result in Section 5.
It uses one integral of the complete Ricci residual, not monotonicity
under symmetrization. In the same ordinary-bolt class let
\(\mu=r_0^2>0\), and choose R in the unchanged CH collar. No bound on
F,w,phi or their derivatives is assumed, apart from positivity and the
smooth bolt and collar conditions. In particular
\(w=(r-r_0)^2\widetilde w(r)\) and phi is smooth in r at the bolt, as in
the smooth profile class of I54. Define \[
 A_0=D_0-2F(w')^2,\qquad A_3=D_3-z_0,\qquad A_H=D_H+c_0,
 \qquad V_w=U-2F\phi'w',
\] \[
 X=2A_H-A_0-A_3,\qquad Y=A_0+A_3+2A_H,
 \qquad Z=A_0-A_3.
\] Here Y is exactly \(e^{2\phi}(s_G+4)\). Put \[
\begin{aligned}
 I_R&=\int_{r_0}^Rr^{-5}\,dr
                  =\tfrac14(\mu^{-2}-R^{-4}),\\
 P_w&=\int_{r_0}^R\left\{
    \frac{2F(w')^2}{r}+\frac{8(\cosh(2w)-1)}{r^3}
                 +\frac{8\sinh^2(2w)}{r^3F}\right\}\,dr,\\
 T_\phi&=2\int_{r_0}^R\frac{F(\phi')^2}{r}\,dr,
 \qquad M=1+2/\mu+P_w-T_\phi .
\end{aligned}                                                     \tag{8.1}
\] These are integrals of the actual profiles. Both P\_w and T\_phi are
nonnegative and finite for smooth bolt data.

\textbf{Proposition 5.} Every such conformal triaxial filling satisfies
\[
\boxed{\begin{aligned}
 \frac{J(G)}{2\pi^2}={}&\frac{M^2}{4I_R}
 +\frac14\int_{r_0}^R r^3
                    \left(X-\frac{M}{I_Rr^4}\right)^2dr\\
 &+\frac14\int_{r_0}^Rr^3Y^2dr
  +\frac12\int_{r_0}^Rr^3Z^2dr
  +2\int_{r_0}^Rr^3V_w^2dr .
\end{aligned}}                                                    \tag{8.2}
\] In particular the whole nonconformal class \(\phi=0\), including
arbitrary finite triaxial anisotropy and arbitrary positive F, obeys \[
 \boxed{\frac{J(K_{F,w})}{2\pi^2}
 \ge\frac{(\mu+2+\mu P_w)^2}{1-\mu^2/R^4}>4.}           \tag{8.3}
\]

\emph{Proof.} Direct substitution in the actual residual gives \[
\begin{aligned}
 X={}&F''+3F'/r-8(F-1)/r^2
       +2(F\phi')'-2F\phi'/r-2F(\phi')^2\\
    &+2F(w')^2+2c_0+z_0 .
\end{aligned}
\] After multiplication by 1/r, its conformal linear terms are exactly
\(2(F\phi'/r)'\). Their boundary values vanish: F = 0 at the bolt, phi
is radially smooth there, and phi = 0 on the full collar. The F terms
have integral \[
 \left[F'/r+4F/r^2\right]_{r_0}^R
                    +8\int_{r_0}^Rr^{-3}dr=1+2/\mu.
\] This uses F'(r\_0) = 2/r\_0 and the exact outer values F(R) = 1+R
squared/6, F'(R) = R/3. Hence \(\int X/r=M\). Pointwise, \[
 A_0^2+A_3^2+2A_H^2=(X^2+Y^2)/4+Z^2/2.
\] Projecting X in \(L^2(r^3dr)\) onto the single function r to the
power minus four proves (8.2). Set T\_phi = 0 for (8.3). QED.

This result excludes the whole diagonal almost-Kahler class itself, not
only small anisotropy or favorable D signs. It does not contradict the
earlier descending plateau: a decrease from one high-cost metric to
another need not cross the threshold four. It also does not prove the
corresponding result for all conformal factors.

More generally put \[
 \delta_{\mu,R}=1+\frac2\mu
                  \left(1-\sqrt{1-\mu^2/R^4}\right)>1.
\] The entire class with \(T_\phi-P_w\le\delta_{\mu,R}\) has J greater
than 8 pi squared. Indeed M is then at least
\(2\mu^{-1}\sqrt{1-\mu^2/R^4}>0\), and (8.2) gives the nonstrict bound;
equality would require X = M/(I\_R r to the fourth) on the full collar,
where the actual residual X vanishes.

Thus every actual subthreshold conformal cell necessarily satisfies \[
 \boxed{1+\frac2\mu\left(1-\sqrt{1-\mu^2/R^4}\right)
 < T_\phi-P_w
 <1+\frac2\mu\left(1+\sqrt{1-\mu^2/R^4}\right).}         \tag{8.4}
\] It must also make all four remaining squares in (8.2) small enough
simultaneously. This is a complete signed integral restriction on actual
metrics, not a formal choice of four independent variables. For a fixed
improvement \(J/(2\pi^2)\le4-\epsilon\), where \(0<\epsilon<4\), the
sharper necessary condition is \[
 \mu(T_\phi-P_w)\ge
        \mu+2-\sqrt{4-\epsilon}\sqrt{1-\mu^2/R^4}.       \tag{8.5}
\] Thus along mu tending to zero a fixed improvement requires a net
weighted conformal gradient budget of order at least mu inverse. Merely
entering the interval (8.4) by a bounded amount does not produce such an
improvement. This is a requirement on complete smooth one-bolt cells; it
is not a theorem decomposing an arbitrary multi-core filling into
separately chargeable one-bolt regions. Controlling that middle interval
from the scalar, radial and anisotropic equations is still open. The
radial identity \[
 Z=-2F\{\phi''-\phi'/r-(\phi')^2+(w')^2\}+z_0
\] is retained rather than treating Z as a freely assignable residual.

\hypertarget{scalar-coupling-and-an-actual-zero-moment-test}{%
\subsubsection{8.1 Scalar coupling and an actual zero-moment
test}\label{scalar-coupling-and-an-actual-zero-moment-test}}

The scalar equation supplies a second integral, but not a bound on
T\_phi alone. Set \[
 G_\phi=\int r^3F(\phi')^2dr,\qquad
 B_\phi=\int r^3(e^{2\phi}-1)dr,\qquad n=N/(2\pi^2).
\] The scalar identity of I54 and the vanishing endpoint term
\([r^3F\phi']=0\) give \[
 L_Y:=\int r^3Ydr=4B_\phi-\mu^2-2\mu-n-6G_\phi,
 \qquad \mu^2T_\phi/2\le G_\phi\le R^4T_\phi/2.         \tag{8.6}
\] Thus the Y-square in (8.2) splits exactly into \[
 \frac{L_Y^2}{R^4-\mu^2}
 +\frac14\int r^3
           \left(Y-\frac{4L_Y}{R^4-\mu^2}\right)^2dr.    \tag{8.7}
\]

The main radial moment M can nevertheless vanish on genuine smooth
metrics. Take the complete cutoff Kahler recovery F\_mu with fixed
physical cutoff radius, w = 0, and a nonzero
\(\eta\in C_c^\infty((2,3))\). For sufficiently small mu put \[
 C_\mu=2\int_2^3F_\mu(\sqrt\mu x)(\eta'(x))^2/x\,dx,
 \qquad a_\mu^2=(\mu+2)/C_\mu,
 \qquad \phi(r)=a_\mu\eta(r/\sqrt\mu).
\] Here C\_mu tends to the positive integral with F replaced by
\(1-x^{-2}\). The metric is smooth and preserves both the bolt and full
CH collar. Direct change of variable gives \[
             P_w=0,\qquad T_\phi=1+2/\mu,\qquad M=0.
\] Its entire cost remains uniformly bounded: the perturbation is
supported on a region of volume O(mu squared), and its curvature is O(mu
inverse) with uniformly controlled rescaled profiles. It is still
conformal to a genuine Kahler cell, so its cost is strictly greater than
8 pi squared by I54. In (8.2) that positive cost is carried by the
remaining squares. In particular neither a bound on M away from zero nor
the universal inequality T\_phi \textless= P\_w + 1 can finish the
conformal problem. The missing estimate must couple the full scalar
variance, Z and V\_w, rather than discard them after taking the two
moments.

\hypertarget{a-two-scale-requirement-for-any-fixed-improvement}{%
\subsection{9. A two-scale requirement for any fixed
improvement}\label{a-two-scale-requirement-for-any-fixed-improvement}}

The radial moment and the actual almost-Kahler transfer restrict where
the remaining conformal payment could occur. This goes beyond choosing
formal values of T\_phi and P\_w.

\textbf{Proposition 6.} Fix R and \(0<\epsilon<4\). Suppose a sequence
of the complete cells in Section 8 has mu tending to zero and
\(j=J/(2\pi^2)\le4-\epsilon\). For every \(0<\alpha<1/4\), \[
 \mu T_{\phi,\{r\ge\mu^\alpha\}}\longrightarrow0,
 \qquad
 \liminf\mu T_{\phi,\{r\le\mu^\alpha\}}
                       \ge2-\sqrt{4-\epsilon}>0.         \tag{9.1}
\] At the same time the complete metric changes cannot be confined to a
CH background region whose volume tends to zero.

\emph{Proof.} Use normalized quantities \[
 d=R^4/4,\qquad v=(R^4-\mu^2)/4,\qquad
 m=R^4+2\mu,\qquad n=N/(2\pi^2).
\] These equal the quantities of I54 divided by 2 pi squared; this
normalization is confined to this section. The conformal transfer
remainder is not discarded: with u = exp(-phi), \[
 \frac{6\int|du|_G^2dV_G}{2\pi^2}
                       =6\int r^3F(\phi')^2dr=6G_\phi.
\] The actual scalar and Weyl bounds of I54 therefore give \[
 j\ge4-4d+
       \frac{(m+n+6G_\phi)^2+2(m-2n)^2}{12v}.
\] The omitted term is \(\int(s_G+4)^2/(8\pi^2)\ge0\). The exact
algebraic decomposition \[
 (m+n+6G)^2+2(m-2n)^2
             =9\left(n-\frac{m-2G}{3}\right)^2+2(m+4G)^2
\] implies \[
 0\le G_\phi\le
             \frac{\sqrt{24dv-6\epsilon v}-m}{4}=O_R(1). \tag{9.2}
\] If the radicand or right side is negative, the hypothesized cell is
already excluded. On the exterior of any canonical radius rho, \[
 T_{\phi,\{r\ge\rho\}}\le2G_\phi/\rho^4.
\] Take rho = mu to the power alpha. The first conclusion in (9.1)
follows because 1 - 4 alpha is positive. Equation (8.5), together with
P\_w nonnegative, supplies the lower bound on mu T\_phi; subtract the
vanishing exterior contribution to obtain the second conclusion.

Finally, if all changes were confined to a CH region of background
volume d\_s tending to zero, the general complete-cell inequality in I54
would give \(J\ge8\pi^2-4d_s/3\), contradicting the fixed improvement.
More precisely it requires \(d_s\ge3\pi^2\epsilon/2\). This last bound
concerns the entire modified background, not the canonical size of the
exceptional bolt. QED.

Consequently a putative fixed improvement in this class needs both a
microscopic conformal payment and a nonvanishing-scale modification of
the full geometry. A small-support conformal bubble alone cannot provide
it; neither can a uniformly bounded weighted conformal gradient on the
shrinking canonical scales. The proposition does not classify the
transition or prove that the two required mechanisms cannot coexist.
That simultaneous compatibility, including all the squares of (8.2), is
the remaining construction test in this class.

\end{document}
